\begin{lemma}[Lemma rutcor: reserved colors and the terminal PE]
Let $(T,\varphi_n)$ be a positive-rung extended Tashkinov tree, let
$T_n$ be its last ladder tree, and let $T_{n,0}^{*}$ be a hierarchy
base with closure tree $R_n$.  Let
\[
 T_n=T_{n,0}\subset T_{n,1}\subset\cdots\subset
 T_{n,q}\subset T_{n,q+1}=T
\]
be a good hierarchy, where $q\geq1$, and put
$T_{n,j}^{*}=\noderef{HierarchyStarredStage}(H,j)$.  Thus
$T_{n,q}^{*}=T_{n,q}$.  Let $p=\noderef{HierarchyPathNumber}(H)$ and
let $T(y_{p-1})$ be the ordered prefix of $T$ of length
\[
 \min\{|T_{n,q}|+(p-1),|T|\}.
\]
Assume the exact prefix-elementarity assertion (6.6), namely that
$V(T(y_{p-1}))$ is elementary under $\varphi_n$.  Also assume that
$V(T_n)$ is elementary under the pre-step coloring
$\varphi_{n-1}$; formally this is the coloring at the predecessor of
the last positive index in the certified series.

Fix $\alpha\in\overline\varphi_n(T_{n,q})$.  For two nested displayed
trees $A\subseteq B$, the expression
$\varphi_n\langle B-A\rangle$ used below means, inline and with exact
edge identities,
\[
 \left\{c\in[k]:\text{ for some }t\in\operatorname{Fin}|B|,
       |A|\leq t\text{ and }\varphi_n(B.\mathit{edge}(t))=c\right\}.
                                                               \tag{*}
\]
If at least one $T_{n,i}^{*}$, $0\leq i\leq q$, is
$\alpha$-closed, then for every maximal such index $r$, equivalently
for every $r$ such that $T_{n,r}^{*}$ is $\alpha$-closed and no
$T_{n,s}^{*}$ with $r<s\leq q$ is $\alpha$-closed,
\[
       \alpha\notin
       \varphi_n\langle T_{n,q}-T_{n,r}^{*}\rangle.            \tag{1}
\]

If no $T_{n,i}^{*}$, $0\leq i\leq q$, is $\alpha$-closed, then all
of the following assertions hold simultaneously.  For every
$1\leq j\leq q$,
\[
 \alpha\in
   \bigcup_{\eta\in D_{n,j}}\Gamma^{j-1}_{\eta}
 \quad\text{and}\quad
 \alpha\in
   \bigcup_{\eta\in D_{n,j-1}}\Gamma^{j-1}_{\eta}
   =:\Gamma^{j-1}.                                             \tag{2}
\]
The terminal extension kind is $\Theta_n=PE$.  Moreover, if
$v_\alpha$ is the vertex of $T$ at
$\noderef{HierarchyFirstMissingPosition}(\varphi_n,T,\alpha)$, then
\[
                  v_\alpha\in V(T_n)\setminus V(R_n),          \tag{3}
\]
and
\[
       \alpha\notin
       \varphi_n\langle T_{n,q}-T_n\rangle.                    \tag{4}
\]
The two unions in (2) are respectively the reserved union indexed by
$D_{n,j}$ and the full paper union $\Gamma^{j-1}$ indexed by
$D_{n,j-1}$.
\end{lemma}
\begin{proof}
We first fix the ordered-stage facts and the elementarity instances that
will be used in both branches.  Since $q>0$, the last starred stage from
\noderef{HierarchyStarredStage} is the unstarred stage $T_{n,q}$.
The nested and terminal fields of \noderef{TashkinovHierarchy} make it
an ordered initial segment of $T$, and every earlier hierarchy stage and
the base core $T_{n,0}^{*}$ is an ordered initial segment of it.

Unfold \noderef{HierarchyPathNumber}.  The final strict hierarchy
inclusion gives $|T_{n,q}|<|T|$.  The path number is either an element of
the nonempty candidate interval beginning at $1$, or the positive tail
length $|T|-|T_{n,q}|$; hence $p\geq1$.  Consequently the prefix
$T(y_{p-1})$ contains $T_{n,q}$ entry by entry.  The exact-prefix
property of \noderef{TreeSequencePrefix} and transitivity of the
initial-segment relations therefore put the vertex set of every
$T_{n,j}$ and every $T_{n,j}^{*}$ inside
$V(T(y_{p-1}))$.  Restricting the pairwise-disjointness condition in the
assumed \noderef{ColoredVertexSetElementary} property proves that all
these vertex sets, including the PE core when it occurs, are elementary
under $\varphi_n$.  This is the promised derivation of stage
elementarity from (6.6), rather than a separate elementarity assumption.

We next record a downward-propagation claim.  Suppose
$1\leq t\leq q$ and no starred stage with index in $[t,q]$ is
$\alpha$-closed.  Then
\[
                    \alpha\in\overline\varphi_n(T_{n,t}).      \tag{5}
\]
We prove this by descending induction on $t$.  At $t=q$, (5) is the
hypothesis that $\alpha$ is missing on $T_{n,q}$.  For the induction
step, suppose that $\alpha$ is missing on $T_{n,t+1}$ but not on
$T_{n,t}^{*}=T_{n,t}$.  Every witness at which it is missing then lies
in
\[
             V(T_{n,t+1})\setminus V(T_{n,t}^{*}).
\]
The elementarity of $T_{n,t+1}$ established above lets us apply
\noderef{GoodHierarchyNewMissingClosed}; that lemma says that a color
newly missing in this difference is closed on $T_{n,t+1}$.  Since
$t+1>0$, this is precisely the starred stage
$T_{n,t+1}^{*}$, contrary to the assumption on $[t,q]$.  Thus the color
was already missing on $T_{n,t}$, proving (5).

Assume first that a closed starred stage exists, and fix an index $r$
which is closed and has no later closed index.  If $r=q$, the index set
in (*) for $T_{n,q}-T_{n,r}^{*}$ is empty, so (1) is immediate.  Suppose
$r<q$, and fix $s$ with $r<s\leq q$.  Maximality of $r$ says that no
starred stage in $[s,q]$ is closed, so (5) gives
$\alpha\in\overline\varphi_n(T_{n,s})$.  On the other hand,
$T_{n,s}=T_{n,s}^{*}$ is not $\alpha$-closed.

Clause `minusClosed' of \noderef{GoodHierarchy}, expanded through
\noderef{TreeColorSetMinusClosed}, says that every color missing on
$T_{n,s}$ but outside
\[
                   C_s:=\bigcup_{\eta\in D_{n,s}}
                                  \Gamma^{s-1}_{\eta}
\]
has empty boundary there.  If $\alpha\notin C_s$, this clause and the
just-proved missingness would say exactly that $T_{n,s}$ is
$\{\alpha\}$-closed in the sense of
\noderef{TreeColorSetClosed}, a contradiction.  Hence there is
$\eta\in D_{n,s}$ with
\[
                         \alpha\in\Gamma^{s-1}_{\eta}.          \tag{6}
\]

We turn (6) into an exact segment exclusion.  By
\noderef{HierarchyStoredColorsAntitone},
$D_{n,s}\subseteq D_{n,s-1}$, so the `avoidance' field of
\noderef{GoodHierarchy} at index $s-1$ applies to this $\eta$ and
$\alpha$.  Also, the definition
\noderef{HierarchyStoredColorsAt} says that
$\eta\in D_{n,s}$ excludes $\eta$ from
$\overline\varphi_n(T_{n,s}^{*})=
  \overline\varphi_n(T_{n,s})$.
Because $T_{n,s}$ is an ordered prefix of $T$, the definition
\noderef{HierarchyFirstMissingPosition} then puts the first
$\eta$-missing position at or after $|T_{n,s}|$ (with the last position
used if there is no missing witness at all).  Thus the minimum occurring
in \noderef{HierarchySegmentColors} is $|T_{n,s}|$.  Its exact
edge-index definition becomes
\[
 \left\{c:\exists t\in\operatorname{Fin}|T_{n,s}|,
   |T_{n,s-1}^{*}|\leq t,
   \varphi_n(T_{n,s}.\mathit{edge}(t))=c\right\}.
\]
The avoidance field and (6) say that $\alpha$ is absent from this set.
In other words,
\[
 \alpha\notin
 \varphi_n\langle T_{n,s}-T_{n,s-1}^{*}\rangle.               \tag{7}
\]

Equation (7) holds for every $s=r+1,\ldots,q$.  The starred stage with
positive index $s-1$ is the preceding unstarred stage, while at
$s=1$ it is the base core.  The ordered-prefix equalities in
\noderef{TashkinovHierarchy} therefore partition the exact half-open
edge-index interval
\[
              [,|T_{n,r}^{*}|,|T_{n,q}|,)
\]
into precisely the intervals used in (7).  Initial-segment equality
identifies each displayed stage edge with the edge of $T_{n,q}$ at the
same index.  Taking the finite union of (7) proves (1).

Now suppose that no starred stage is $\alpha$-closed.  Applying (5) for
every $1\leq j\leq q$, and repeating the minus-closed argument above,
gives
\[
        \alpha\in\bigcup_{\eta\in D_{n,j}}
                            \Gamma^{j-1}_{\eta}.                \tag{8}
\]
For a witness $\eta\in D_{n,j}$ in (8),
\noderef{HierarchyStoredColorsAntitone} gives
$\eta\in D_{n,j-1}$.  The same color belongs to the same pair
$\Gamma^{j-1}_{\eta}$, so (8) also gives
\[
        \alpha\in\bigcup_{\eta\in D_{n,j-1}}
                            \Gamma^{j-1}_{\eta}=\Gamma^{j-1}.
\]
This proves both memberships in (2), including the larger paper union.

At $j=1$, choose $\eta\in D_{n,1}$ with
$\alpha\in\Gamma^0_\eta$.  Antitonicity gives
$\eta\in D_{n,0}$.  The avoidance field of
\noderef{GoodHierarchy} at index $0$ puts every member of
$\Gamma^0_\eta$ in
\noderef{HierarchyStageMissingColors} at index $0$, which by definition
is $\overline\varphi_n(T_n)$.  Therefore
\[
                         \alpha\in\overline\varphi_n(T_n).     \tag{9}
\]

We prove that the last extension is PE by inspecting its actual
certificate.  The positive-rung data in
\noderef{ExtendedTashkinovTree} identifies $\varphi_n$ with the last
coloring in \noderef{TashkinovSeries}; its positive-index certificate is
one of the SE, PE, or RE constructors in
\noderef{TashkinovSeriesStep}.  Suppose it is not PE.

In the RE constructor, the current coloring is the pre-step coloring,
and the constructor says that $V(T_n)$ is closed under that coloring.
Together with (9), the definition
\noderef{ColoredVertexSetClosed} makes the boundary of $\alpha$ at
$T_n$ empty.  In the SE constructor, write $\pi$ for the selected current
coloring.  The `previousClosed' field closes $T_n$ under the pre-step
coloring, while the `selectedStable' field of
\noderef{MaximumDefectiveChoice} says that $\pi$ is stable on $T_n$.
Both colorings are proper and leave exactly $e_0$ uncolored.  Symmetry
from \noderef{sc2} therefore gives the reverse stable relation.  If a
boundary identity had color $c$ under $\pi$, with $c$ missing at a tree
vertex under $\pi$, then \noderef{BoundaryEdges} and
\noderef{TreeIncidentEdge} say that this identity is incident with
$T_n$.  Reverse stability would give the same edge value under the
pre-step coloring, and the vertex clause of
\noderef{StableColoring} would make $c$ missing at the same vertex
there.  This contradicts `previousClosed'.  Thus $T_n$ is closed under
$\pi=\varphi_n$ in the SE case as well.  In either non-PE case, (9)
makes $T_n$ $\alpha$-closed.  The nonparallel field of
\noderef{HierarchyBase} identifies $T_{n,0}^{*}$ with $T_n$, contrary
to the no-closed-stage hypothesis.  Hence
\[
                              \Theta_n=PE.                       \tag{10}
\]

Let $a=\noderef{HierarchyFirstMissingPosition}
(\varphi_n,T,\alpha)$ and $v_\alpha=T.\mathit{vertex}(a)$.  Since
$T_n$ is an ordered initial segment of $T$, (9) and the least-candidate
definition of $a$ show that $v_\alpha$ is an actual missing-color witness
and
\[
                              v_\alpha\in V(T_n).                \tag{11}
\]
We now use the pre-step elementarity premise, and only here use the
special PE closure theorem.  Apply
\noderef{ParallelCommonMissingClosed} with (10), the generating-coloring
identity from \noderef{ExtendedTashkinovTree}, and the stated
elementarity of $V(T_n)$ under $\varphi_{n-1}$.  It says that the PE core
$T_n\vee R_n$ is closed for every color missing on both $T_n$ and
$R_n$.  If $v_\alpha$ also belonged to $V(R_n)$, then it would itself
witness that $\alpha$ is missing on $R_n$; equation (9) would make
$\alpha$ common-missing.  The cited lemma would make
$T_{n,0}^{*}=T_n\vee R_n$ $\alpha$-closed, contradicting the
no-closed-stage hypothesis at index $0$.  Therefore (11) strengthens to
\[
                        v_\alpha\in V(T_n)\setminus V(R_n),
\]
which is (3).  Notice that the PE certificate alone was not used to
deduce common-missing closure: the exact lower-rung elementarity premise
of \noderef{ParallelCommonMissingClosed} was supplied.

For completeness, we next derive the two PE-base exclusions used in the
last edge argument.  The parallel field of
\noderef{HierarchyBase} and \noderef{OrderedTreeUnion} identify the
core vertex set with $V(T_n)\cup V(R_n)$.  Its elementarity was already
obtained by restricting (6.6).  Since $v_\alpha$ lies in this core and
misses $\alpha$, no different core vertex can miss $\alpha$.  In
particular,
\[
 \alpha\notin
 \bigcup_{w\in V(R_n)\setminus V(T_n)}
          \overline\varphi_n(w).                                \tag{12}
\]

The same elementarity excludes $\alpha$ from every exact edge of
$R_n$ not already in $T_n$.  Indeed, the PE field of
\noderef{HierarchyBase} supplies a
\noderef{TashkinovClosure} of the exact support prefix $T_n(v_n)$ whose
terminal tree is $R_n$.  By
\noderef{TashkinovAugmentationSequence}, this is a finite sequence of
steps from \noderef{TashkinovAugmentation}.  In one such step the unique
new last edge has a color missing on the preceding tree.  Every preceding
tree is an initial segment of $R_n$, so a hypothetical new edge of color
$\alpha$ would supply a vertex $w\in V(R_n)$ missing $\alpha$.
Equation (3) gives $v_\alpha\notin V(R_n)$, hence $w\ne v_\alpha$;
both vertices lie in the elementary core, contradicting that both miss
$\alpha$.  Induction along the finite augmentation sequence excludes
$\alpha$ from every edge added after the support prefix.  Every edge of
that prefix is the same exact edge identity as an edge of $T_n$, by
\noderef{TreeSequencePrefix}.  Consequently
\[
 \text{no identity in }E(R_n)\setminus E(T_n)
       \text{ has color }\alpha,                               \tag{13}
\]
where the exact edge sets are those of
\noderef{TreeSequenceEdgeSet}.
This proves the missing-color exclusion (12) and the edge-color exclusion
(13) directly from (6.6) and the Tashkinov augmentation data.

It remains to pass from the PE base to $T_{n,q}$.  For every
$1\leq s\leq q$, the witness in (8), antitonicity, avoidance, and the
exact first-missing-position calculation used for (7) give
\[
 \alpha\notin
 \varphi_n\langle T_{n,s}-T_{n,s-1}^{*}\rangle.                \tag{14}
\]
Uniting the adjacent half-open index intervals in (14) gives
\[
 \alpha\notin
 \varphi_n\langle T_{n,q}-T_{n,0}^{*}\rangle.                 \tag{15}
\]
Finally, the `leftSegment' and `newEntryFromRight' fields of
\noderef{OrderedTreeUnion} say that every exact entry of
$T_{n,0}^{*}=T_n\vee R_n$ after $T_n$ is an entry of $R_n$ with the
same edge identity.  Injectivity of tree-sequence edge entries ensures
that this right-hand identity is not already an edge of $T_n$.
Therefore (13) excludes $\alpha$ from the entire interval between
$T_n$ and the PE base, while (15) excludes it from the interval after
the base.  These two adjacent exact index intervals form
$[|T_n|,|T_{n,q}|)$, so their union is precisely (*) for
$T_{n,q}-T_n$.  This proves (4) and completes the second branch.
\end{proof}
